\documentclass[11pt]{article}
\usepackage[margin=1in]{geometry}
\usepackage[dvipsnames]{xcolor}
\definecolor{darkgray}{gray}{0.35}
\usepackage{enumitem}
\usepackage[round]{natbib}
\usepackage{hyperref}
\usepackage{xr-hyper}
\externaldocument{manuscript}
\externaldocument{appendix}
\usepackage{soul}
\usepackage{graphicx}
\usepackage{float}

\newcommand{\comment}[1]{\textbf{Comment:} \textit{#1}}
\newcommand{\response}[1]{\textbf{Response:} #1}
\sethlcolor{yellow}
\newcommand{\placeholder}[1]{\hl{#1}}
\newcommand{\todo}{\textcolor{red}{\textbf{[TODO]}}}
\newcommand{\manuscript}[2]{\begin{quote}\placeholder{#1:}\\{\color{darkgray}``#2''}\end{quote}}

\begin{document}

\noindent\textbf{\Large Response to Reviewers}\\[0.5em]
\noindent Manuscript: PMEDICINE-D-25-04465R1\\
\noindent Title: Reducing spatial clustering to prevent \emph{Mycobacterium tuberculosis} transmission in a busy Zambian hospital: A modelling study based on person movements, environmental and clinical data\\[1em]

\noindent Dear Dr.\ Farrell,\\[0.5em]

\noindent Thank you for the opportunity to revise our manuscript. We are grateful for the constructive feedback from all three reviewers, which has substantially improved the paper. Below, we provide point-by-point responses to all editorial and reviewer comments. Changes in the manuscript are indicated by page and line numbers, and copies of all revised and newly added figures are appended at the end of this letter.\\[0.5em]

\noindent The most important changes in the revised manuscript are:

\begin{enumerate}
\item \textbf{Clearer distinction between empirical and modelled findings.} Throughout the manuscript, we now consistently qualify transmission risk estimates as modelled outputs and use attenuated language (``potentially prevented,'' ``may have a meaningful public health impact'') rather than definitive claims (\ref{r:2.1}, \ref{r:2.7}).

\item \textbf{Expanded description of the study setting and interventions.} The Methods section has been restructured into separate ``Study setting'' and ``Study interventions'' subsections, with the intervention details (former Supplementary Table S1) moved into the main text as Table 1. A new Figure~\ref{fig:interventions} illustrates the intended impact of the interventions during each study phase (\ref{r:2.12}, \ref{r:2.13}, \ref{r:2.18}).

\item \textbf{Additional sensitivity analyses.} We added three new sensitivity analyses in response to reviewer concerns: (i) an alternative kernel (exponential) and metric (Shannon entropy) for measuring spatial clustering (\ref{r:3.3}, \ref{r:2.5}); (ii) a revised HIV sensitivity scenario assuming 33\% lower infectiousness, replacing the implausible six-fold increase (\ref{r:2.31}, \ref{r:3.4}); and (iii) a scenario in which all confirmed and presumptive TB patients are considered infectious (\ref{r:2.8}).

\item \textbf{Discussion of modelling limitations.} A new paragraph discusses the limitations of the well-mixed assumption, the pros and cons of computational fluid dynamics vs.\ Wells-Riley-based approaches, and the need for empirical airflow data (\ref{r:2.27}, \ref{r:2.3}).

\item \textbf{Reporting improvements.} Absolute ventilation rates (L/s) are now reported alongside air exchange rates (\ref{r:2.14}). Prevented infections are reported as daily rates and per-entry rates (\ref{r:2.9}, \ref{r:2.15}). The colour scheme in figures showing intervention effects and weekday-specific data has been updated (\ref{r:2.19}, \ref{r:2.36}).
\end{enumerate}

\clearpage

% ============================================================
\section*{Editorial Comments}
% ============================================================

\begin{enumerate}[label=\textbf{E\arabic*.}, leftmargin=*]

\item\label{r:E1} \comment{Please clarify the Data Availability Statement: where will the data be deposited?}

\textcolor{red}{\response{We have revised and clarified the Data Availability Statement. Environmental and person-tracking data have been made publicly available. Clinical data have not been deposited, as the individual-level records pose an unnecessary risk of inadvertent patient disclosure (see also \ref{r:2.24}). However, we provide the full code to replicate our simulation results, which includes preprocessed datasets in which person tracks from TB patients are indicated---but without any further personal attributes such as age, sex, or HIV status.}}

\item\label{r:E2} \comment{Please provide a doi for the code.}

\response{We have deposited the code and preprocessed data (see \ref{r:E1}) in a public repository and provide a DOI: \placeholder{DOI link}.}

\item\label{r:E3} \comment{Please add funder URLs to the manuscript metadata.}

\textcolor{red}{\response{We have added the funder URLs to the manuscript metadata.}}

\item\label{r:E4} \comment{Please add a statistics subsection to the Methods.}

\response{We have renamed the existing Methods subsection ``Estimation of intervention effects on \emph{Mtb} transmission'' to ``Statistical analysis: Estimation of intervention effects on \emph{Mtb} transmission'' (p.\,\pageref{chg:E4}). This subsection describes the Bayesian regression model used to estimate intervention effects, including the adjustment for weekday effects and person-time, and thus serves as the statistical analysis section.}

\end{enumerate}

\clearpage

% ============================================================
\section*{Reviewer \#1}
% ============================================================

\noindent\textit{``Thank you for the opportunity to review this manuscript. As the authors say, a simple, low-cost operation can make a difference in TB transmission and they have provided evidence to support this in a busy hospital setting in Africa. Overall, the manuscript is well written.''}

\medskip
\noindent We thank Reviewer \#1 for the positive assessment and the constructive questions. In response, we have clarified the rationale for the study duration and baseline design, explained the nature of the technical issues, and replaced the term ``incident TB'' with ``detected TB cases.'' All changes are detailed below.

\begin{enumerate}[label=\textbf{R1.\arabic*.}, leftmargin=*]

\item\label{r:1.1} \comment{In the methods you say you had a sample of 14 days to capture sufficient variation. Is this sufficient time and is there any reference for this?}

\response{Thank you for raising this point. The 14-day duration per intervention phase was chosen pragmatically, balancing two competing considerations: maximising the number of study days to capture sufficient variation in key drivers of transmission risk (patient attendance, ventilation conditions, and TB burden), while managing the considerable operational burden on local study staff who were responsible for collecting data as well as implementing and monitoring the interventions on-site.

We believe that 14 days per intervention phase provided an adequate observation window for several reasons. First, two full weeks encompass all days of the week twice, which is important because patient attendance patterns at this facility varied substantially between weekdays and weekends (e.g., Mondays and Tuesdays were the busiest days). This allowed our statistical models to adjust for day-of-week effects while still retaining sufficient data points per day type. Second, our Bayesian modelling framework explicitly accounts for multiple sources of variability---including daily fluctuations in occupancy, ventilation, and TB case burden---through Monte Carlo simulations that propagate parameter uncertainty into the final estimates. The resulting credible intervals reflect the uncertainty associated with the observed sample size. Third, by including both pre-intervention (9 days) and post-intervention (15 days) baseline phases, we were able to capture and adjust for temporal trends in routine clinic conditions over the full 52-day study period, further strengthening the reliability of the intervention effect estimates.

We acknowledge that there is no established reference for the required sample size in studies of this nature. Unlike clinical trials with well-defined primary endpoints and established effect sizes, formal sample size calculations for complex, multi-component simulation studies such as ours---which integrate movement tracking, environmental monitoring, and transmission modelling---would require assumptions about numerous interacting parameters, making such calculations somewhat arbitrary and potentially conveying a false sense of precision about the amount of data needed. Instead, the adequacy of the data is better judged by the precision of the resulting estimates: the 95\% credible intervals for the intervention effects on spatial clustering (e.g., 13--32\% reduction for the first intervention) and modelled Mtb transmission risk (e.g., 29--48\% reduction) are sufficiently narrow to support meaningful conclusions.

We have revised the manuscript to better explain the rationale behind the chosen study duration (p.\,\pageref{chg:R1.1}).}

\item\label{r:1.2} \comment{You had 15 days `baseline' at the end ``to capture temporal trends'' but surely processes changed due to the intervention? I am not entirely sure of the purpose of this and would appreciate further explanation in the manuscript.}

\response{Thank you for this question. The post-intervention baseline period served both a practical and an analytical purpose, which we now explain more clearly in the revised manuscript (p.\,\pageref{chg:R1.2}).

\textit{Practical rationale.} The pre-intervention baseline period was shorter than originally planned because several study days had to be excluded due to electricity or internet outages at the facility. These outages interrupted the optical tracking system and reset the internal study ID counter, rendering the affected days unusable for analysis. Because these exclusions were only identified after the intervention phase had already commenced, we compensated by adding 15 days of baseline observations (without any interventions) after the two intervention phases. This ensured that we had a sufficient number of non-intervention days to characterise routine clinic conditions. We note that these technical issues are distinct from the server-side configuration issue mentioned in our response to comment \ref{r:1.3}, which affected only the collection of person-level attributes. 

\textit{Analytical rationale.} This design choice turned out to be additionally beneficial because of seasonal patterns observed over the study period: patient attendance declined steadily from June to August, coinciding with rising ambient temperatures and reduced circulation of respiratory infections. By bookending the intervention phases with baseline periods at the start (higher attendance) and end (lower attendance) of the study, we could verify that our statistical adjustments for temporal variation were adequate. The Bayesian regression model already adjusted for variation in attendance (person-time) and day-of-week effects, so that temporal trends in occupancy did not confound the estimated intervention effects on Mtb transmission risk. In that sense, the post-intervention baseline was not strictly necessary to correct for seasonal patterns, but it provides additional reassurance that the adjusted transmission risk estimates and intervention effects were not driven by temporal trends.

\textit{Potential carryover effects.} We appreciate the reviewer's concern that processes may have changed following the interventions. We believe that carryover effects were limited for several reasons. First, all infrastructure changes (repositioned benches, signage, seat markings) were temporary and removed at the end of the intervention phases. Second, while study staff could theoretically have continued monitoring spatial clustering or advising patients to physically distance, they were aware of the study design and instructed not to continue intervention-related activities during baseline periods. Third, behavioural changes among clinic attendees are unlikely to have persisted, as few patients visit the facility multiple times within the relatively short study period. Consistent with these considerations, the modelled Mtb transmission risk adjusted for person-time was higher during the post-intervention baseline compared to the intervention phases, suggesting that the beneficial effects of the interventions diminished once they were removed (Supplementary Figures \ref{fig:S9} and \ref{fig:S10}).}

\item\label{r:1.3} \comment{You say there were technical issues in first 10 days which covers all your baseline period -- how reliable do you therefore think your baseline data are? I don't recall this being mentioned in the study limitations.}

\response{We appreciate this concern and recognise that this point was unclear in the original manuscript. The technical issues during the first 10 days did not affect the core tracking data (person positions and movements) used in the main analysis. The optical sensors operated without issues throughout the entire study period, continuously recording the location and movement of individuals in the waiting area.

The technical problem was a server-side configuration issue that prevented the collection of person-level attributes, specifically gender and healthcare worker status. These attributes are only used in supplementary subgroup analyses (e.g., comparing transmission risk by gender or between healthcare workers and other attendees) and are not relevant to the main analysis of intervention effects on spatial clustering and Mtb transmission risk. This is why the issue was not discussed in the study limitations.

We have revised the manuscript to clarify the nature of these technical issues and to make explicit that the core tracking data were not impacted (p.\,\pageref{chg:R1.3}).}

\item\label{r:1.4} \comment{Table 1: Use of the term ``incident TB''? I see that you have an explanation of this term as a footnote of the table however it is somewhat misleading as the majority are presumptive cases and not actually new TB cases as is implied with the meaning of the word `incident'. Suggest you use a different word.}

\response{We agree that ``incident TB'' could be misleading, as it implies confirmed new cases in standard epidemiological usage. We have replaced it with ``detected TB cases'' throughout the manuscript, which more accurately reflects that this category includes both presumptive and confirmed cases identified during the study period. The table footnote clarifying the distinction between presumptive and confirmed cases has been retained (Table 1, p.\,\pageref{chg:R1.4a}; table footnote, p.\,\pageref{chg:R1.4b}).}

\end{enumerate}

\clearpage

% ============================================================
\section*{Reviewer \#2}
% ============================================================

\noindent\textit{``Thanks for the opportunity to review this interesting manuscript. I comment as a clinician and an epidemiologist with an interest in Mtb transmission. I would advise that someone with expertise in mathematical modelling be also asked to review this paper.''}

\medskip
\noindent We thank Reviewer \#2 for the thorough and insightful review. The detailed comments have substantially improved the manuscript. In response, we have: restructured the Methods section with expanded descriptions of the study setting, interventions, and environmental monitoring; moved the interventions table into the main manuscript; added absolute ventilation rates; clarified the distinction between empiric and modelled estimates throughout; added a Discussion paragraph on airflow limitations and CFD alternatives; attenuated language regarding the significance of findings; ensured consistent use of TB (disease) vs.\ \emph{Mtb} (pathogen); and addressed numerous additional points as detailed below.

\subsection*{Major Comments}

\begin{enumerate}[label=\textbf{R2.\arabic*.}, leftmargin=*]

\item\label{r:2.1} \comment{Throughout the manuscript -- both when describing study results and when describing other modelling work (e.g.\ Beckwith et al) -- the authors need to make a clear distinction between empiric and modelled estimates of Mtb transmission. E.g.\ the authors have no empiric data to support the assertion (abstract) that the interventions `lowered Mtb transmission risk by an estimated 39\%\dots'}

\response{We agree that the distinction between empiric and modelled estimates must be clear throughout. We have revised the manuscript to consistently qualify transmission risk estimates as modelled rather than empirically observed. Specific changes include:

\begin{itemize}
\item In the Abstract, we now write ``the model estimated that [the interventions] potentially prevented\ldots'' rather than stating that the interventions ``lowered Mtb transmission risk,'' and we distinguish between observed reductions in spatial clustering and model-estimated reductions in transmission risk (p.\,\pageref{chg:R2.1abs2}).
\item In the Results, infection counts and rates are identified as modelled outputs (p.\,\pageref{chg:R2.1results}), while the Gini coefficient changes are presented as empirical observations.
\item In the Discussion, we distinguish between empirical evidence (a more even spatial distribution) and modelled findings (reduced transmission risk) (p.\,\pageref{chg:R2.1disc2},\,\pageref{chg:R2.1disc3}), and have attenuated ``significant public health impact'' to ``may have a meaningful public health impact'' (p.\,\pageref{chg:R2.1disc4}).
\item Throughout, we have replaced definitive language with ``suggest,'' ``potentially,'' and ``estimated.''
\end{itemize}

These revisions make clear that the transmission-related findings are modelled estimates of potential benefits, not empirical evidence of prevented infections.}

\item\label{r:2.2} \comment{The most impressive bit of the paper is the detailed measurement and modelling of person movements. Presumably this approach to measurement has been validated separately. It would be good to include a few lines noting what it was validated against (e.g.\ direct observation) and how well it performed.}

\response{Thank you for this comment. The person tracking system used in our study is an established commercial product developed by XOVIS, a company specialising in optical people-counting and tracking technology. XOVIS sensors are widely deployed in airports, retail environments, and public transport hubs, and the system has been validated and refined through numerous real-world applications. As a commercially verified industry product, we did not conduct a separate validation study against direct observation.

The system was set up and configured for our study setting in collaboration with ASE AG, a Swiss technology company that provides professional installation and configuration services for XOVIS tracking systems. ASE AG's technical staff worked with us to optimise the sensor configuration for the specific layout and conditions of the hospital waiting area, ensuring accurate tracking performance. With the finalised setup, approximately 50\% of persons were tracked end-to-end---meaning their trajectory was recorded uninterrupted from entry to exit of the sensor-covered area, without the track terminating prematurely within the waiting room. ASE AG confirmed that this proportion of end-to-end tracking is within the normal range based on their experience across comparable deployments.

Interrupted tracks do not compromise the main analysis. For quantifying Mtb transmission risk, we averaged over all person tracks on a daily basis to estimate the spatial distribution of individuals in the waiting area; whether individual tracks are complete or fragmented does not affect these aggregate estimates. The only scenario where track interruptions are problematic is for characterising the movement patterns of specific individuals with TB. For these cases, we developed an R Shiny application to manually re-link interrupted track segments belonging to the same individual, thereby reconstructing representative movement trajectories for TB patients. This procedure is described in detail in Supplementary Text A.

We have added a brief description of the XOVIS tracking system and the professional setup support provided by ASE AG to the revised manuscript (p.\,\pageref{chg:R2.2}).}

\item\label{r:2.3} \comment{Whether an intervention that seeks to separate people within a waiting room impacts Mtb transmission is contingent on how well mixed the airspace is. The authors present no data on this. If there are no such data, this should be acknowledged. It would be good if the authors could comment on how well mixed the airspace is likely to be. Are there paddle fans? Are doors being opened and closed? Perhaps there some external data on similar spaces. Do we expect big differences between facilities in air mixing? If so, this will limit the generalisability of the authors conclusions.}

\response{We agree that airspace mixing is an important factor influencing both the overall and individual-level Mtb transmission risk, and we acknowledge that we did not collect empirical data on airflow patterns in the waiting area. We have described the ventilation conditions in the revised manuscript (p.\,\pageref{chg:R2.3vent}): the waiting area was naturally ventilated through the main entrance door and openable windows, with airflow occasionally supported by ceiling fans (paddle fans, as the reviewer asks). The main entrance door was typically kept open during operating hours; windows were opened and closed depending on weather conditions and staff preference. The hospital is situated in a residential area without nearby combustion sources, and replacement air was predominantly outdoor air. We have also added a paragraph in the Discussion to acknowledge and discuss the limitations of the well-mixed assumption and the potential of alternative approaches such as computational fluid dynamics modelling (see also our response to \ref{r:2.27}).

To clarify the modelling approach: our model is an extension of the Wells-Riley framework that incorporates person movements to reflect close-proximity contacts in the estimated Mtb transmission risk. The classical Wells-Riley model assumes a well-mixed airspace, and this assumption is retained in our extension. What our model adds is the ability to capture spatially heterogeneous exposure arising from the proximity of individuals to infectious sources---i.e., it accounts for the fact that people seated closer to a person with TB likely receive a higher dose of infectious quanta than those seated further away. However, the model does not simulate directional airflow or localised air currents that could redistribute quanta in ways that deviate from the well-mixed assumption.

Experimentally determining airflow patterns is possible but challenging, particularly in a longitudinal study such as ours. Airflow in the waiting area is influenced by factors that change continuously---outdoor wind conditions, which doors are open at any given time, the number and movement of people, and temperature differentials---making point-in-time measurements of limited value for a study spanning 52 days. Computational fluid dynamics (CFD) models can simulate localised airflow patterns in greater detail, but they are computationally expensive and typically applied to static or short-duration scenarios, making them difficult to scale to the volume of longitudinal data in our study. We have added a dedicated paragraph to the Discussion addressing the limitations of not measuring airflow, the well-mixed assumption, and the relative strengths and limitations of Wells-Riley-based and CFD-based approaches (p.\,\pageref{chg:R2.27}; see also our detailed response to \ref{r:2.27}).

That said, we believe the well-mixed assumption does not invalidate the estimated intervention effects for two reasons. First, the intervention effects are computed in relative terms---as the ratio of modelled Mtb transmission risk between intervention and baseline phases---rather than as absolute infection counts. While the well-mixed assumption may affect the precision of absolute risk estimates, relative comparisons are more robust, as systematic biases from ignoring airflow patterns would affect both the numerator and denominator similarly. Second, the purpose of our study is to assess and quantify the \textit{potential} of crowding-reduction interventions to lower Mtb transmission risk using a modelling framework specifically designed to incorporate spatial proximity. It is not an empirical study that directly measures prevented infections. We have revised the manuscript to make this distinction clearer and to temper claims accordingly.

Regarding generalisability: we agree that air mixing conditions are likely to differ between facilities---depending on room geometry, ventilation infrastructure, and occupancy patterns---and that this limits the direct transferability of our absolute risk estimates. However, the qualitative finding that reducing spatial clustering can lower proximity-dependent Mtb exposure should hold across settings, even if the magnitude of the effect will vary. We have acknowledged this as a limitation in the Discussion (p.\,\pageref{chg:R2.27}).}

\item\label{r:2.4} \comment{In presenting data on clustering, the authors adjust for day of the week. Looking at Fig S5, it seems that a major determinant of clustering is the overall occupancy of the waiting room. Presumably, this is because, in a full waiting room, all the chairs are occupied. Is day of the week simply a proxy for occupancy and, if so, why not use e.g.\ daily attendance, which will also account for e.g.\ public holidays, strikes, etc? Also, why adjust for day of the week/attendance/occupancy rather than stratify by day of the week/attendance/occupancy? I would expect the intervention to be more effective in a part full than in a crowded waiting room.}

\response{Thank you for this insightful comment. It is important to distinguish between the two outcome measures in our study, as they relate differently to attendance.

\textit{Spatial clustering (Gini coefficient).} The Gini coefficient measures how uniformly people are dispersed across the available space. As a measure of relative spatial inequality, it is inherently normalised for the number of people present: higher attendance means more people per cell, but the Gini coefficient only changes if their spatial distribution becomes more or less uniform. Thus, the Gini coefficient does not mechanically increase with occupancy, and adjusting for attendance is not necessary for this outcome. Day-of-week indicators are included as covariates to capture residual variation in clinic operations (e.g., staffing patterns, types of appointments) that may affect how people distribute themselves in the waiting area. We have elaborated on the Gini coefficient in the revised manuscript (p.\,\pageref{chg:R2.5}).

\textit{Mtb transmission risk.} For the modelled transmission risk, total daily person-time in the waiting area is included as a model offset, effectively normalising the estimated risk by attendance. This accounts for the fact that, all else being equal, more people in the room increases the absolute number of transmission events. Day-of-week indicators are included as additional covariates. 

Regarding stratification by occupancy: we agree that it would be informative to assess whether the intervention effect varies with attendance levels, as one might expect the interventions to be more effective in a partially full than in a crowded waiting room. However, a sample size of 52 study days with only 14 days per intervention phase is insufficient to support such a stratified analysis with adequate precision. This limitation is compounded by the sequential design of the study phases and the seasonal decline in attendance over time, which means that attendance levels are partly confounded with study phase. A study designed to evaluate interaction effects between interventions and attendance would require a larger sample size and, ideally, greater variation in attendance within each study period. We have added this point to the limitations section of the revised manuscript (p.\,\pageref{chg:R2.4}), noting that the sequential design and limited sample size precluded stratification by occupancy, and that future studies with larger sample sizes and greater within-phase variation in attendance could assess potential interaction effects.}

\item\label{r:2.5} \comment{What does a Gini coefficient of 0.55 vs 0.51 mean in practice? Is this a meaningful difference? Can the authors also present some more intuitive metrics, e.g.\ proportion of people sat within 1m of another person, ideally stratified by waiting room occupancy.}

\response{We acknowledge that the absolute Gini coefficient values are not easily interpretable in isolation. To provide more intuition: the Gini coefficient quantifies inequality in how people are distributed across the available space, where 0 represents a perfectly uniform distribution and 1 represents maximal clustering. A value of 0.55 during baseline means that the spatial distribution of people deviated substantially from uniform---people were concentrated in certain areas of the waiting room (typically the front rows near the reception). The reduction to 0.51 during the first intervention indicates that people were more evenly spread across the space, consistent with the redesigned seating layout and distancing measures. While the absolute difference may appear modest, the relative reduction of 24\% (95\% CrI 13--32\%) in the Gini coefficient is substantial and statistically credible.

The Gini coefficient is a well-known measure of inequality with applications in spatial contexts, and it captures precisely what we aim to measure: how uniformly people are dispersed, independent of how many people are present. We have elaborated on this in the revised manuscript (p.\,\pageref{chg:R2.5}).

The alternative metric suggested by the reviewer---the proportion of people seated within 1m of another person---was considered during the statistical analysis but ultimately dropped because it is mechanistically linked to the number of people in the room. Higher attendance naturally increases the expected number of people within 1m distance, regardless of how well they are dispersed. This metric therefore cannot be descriptively compared across study phases without adjusting for attendance. A stratified presentation by attendance level would conflate intervention and attendance effects, and a full cross-stratification by both attendance and intervention phase is not feasible given the limited sample size (see our response to \ref{r:2.4}).

For these reasons, we retained the Gini coefficient as our primary measure of spatial clustering. That said, to ensure our findings are not sensitive to this choice, we have added a sensitivity analysis using both an alternative smoothing method (an exponential kernel, which weights nearby people more heavily than those further away) and an alternative measure of spatial evenness (Shannon entropy, which is perhaps more intuitive: higher values mean people are more evenly spread out). The results were qualitatively consistent across all combinations (Supplementary Table \ref{tab:S3}; see also our response to \ref{r:3.3}).}

\item\label{r:2.6} \comment{As I understand it, the authors sought to model the impact of the intervention by contrasting the outputs of models that assume full mixing vs one that was `proximity based', with a ratio of the number of infections predicted by each model then divided by total person-time. As per comment \#4, I find dividing by person-time problematic, as the impact of the intervention will be different in more vs less crowded waiting rooms. Why not stratify and show that the intervention was modelled to be most impactful when the waiting room was not full?}

\response{Thank you for returning to this important point. We believe it is helpful to distinguish between two related but separate questions: (i) \textit{adjustment}---does the analysis account for the direct influence of attendance on Mtb transmission risk? and (ii) \textit{interaction}---does the intervention effect itself vary with attendance levels?

Regarding adjustment: the rationale for dividing by person-time arose directly from the data. When we initially compared the modelled Mtb transmission risk across study phases without adjustment, we observed a near-monotonic decline over time that closely tracked the seasonal decline in attendance. Dividing by total person-time removed this trend, confirming that the unadjusted decline was driven by attendance rather than by intervention effects. The person-time adjustment thus serves a clear and necessary purpose: it accounts for the fact that higher attendance mechanically increases the absolute number of transmission events through greater spatial density, which is likely the main concern underlying the reviewer's comment. The resulting intervention effects are expressed in relative terms and relative to the number of people attending, so they already factor in that the absolute risk is lower when fewer people visit the clinic.

Regarding interaction: we think the reviewer's suggestion to stratify by occupancy implicitly asks a different question---not whether attendance \textit{confounds} the intervention effect (which is what the person-time adjustment addresses), but whether attendance \textit{modifies} the intervention effect. In other words, is the relative reduction in transmission risk per person-unit-time itself larger or smaller depending on how crowded the room is? For example, it is plausible that on quiet days, distancing measures are easier to enforce and more effective, whereas on busy days the interventions may have less impact because there is simply not enough space to redistribute people. This would constitute an interaction between occupancy and intervention effectiveness---a fundamentally different concern from the confounding bias that the person-time adjustment removes. Our current analysis estimates an average intervention effect across all observed occupancy levels; it does not estimate whether this effect varies with occupancy. Assessing such an interaction would require a stratified or interaction-term analysis, which, as discussed in our response to \ref{r:2.4}, is not feasible given the limited sample size of 14 days per intervention phase and the confounding of attendance levels with study phase due to the sequential design and seasonal trends. We have acknowledged this limitation in the revised manuscript (p.\,\pageref{chg:R2.4}; see also \ref{r:2.4}).

In summary, dividing by person-time is not problematic but rather essential for isolating the intervention effect from the dominant influence of attendance on absolute transmission risk. The question of whether interventions are differentially effective at different occupancy levels remains open and is acknowledged as a limitation requiring further study (p.\,\pageref{chg:R2.4}).}

\item\label{r:2.7} \comment{I believe that, in the proximity based model, quanta were assumed to diffuse radially and there was no mixing. In reality, there will be some mixing. The true situation will likely fall somewhere between the two scenarios the authors modelled. That the `model estimated that short-range exposure remained an important driver of transmission' is only because of the assumptions made about air mixing. Given we have no data on the true extent of mixing, are the authors able to model the number of infections predicted/averted under various assumptions about mixing, e.g.\ full mixing (no intervention effect), some mixing, diffusion only? The authors could then comment on the need for empiric data to tell us which scenario best reflects reality. Ditto, assumptions about viability of bioaerosols, settling (which will depend on the extent of mixing), etc, where there is also little empiric data to guide model assumptions. The claim that the authors' findings present `a significant public health impact in reducing the risk of Mtb transmission in clinical settings' seems overblown.}

\response{We appreciate these thoughtful and important points. The reviewer is correct that the true situation likely falls somewhere between the two scenarios we modelled (full radial diffusion vs.\ well-mixed airspace), and that the claim of short-range exposure as a driver of transmission is contingent on the modelling assumptions about air mixing. We have added a paragraph to the Discussion that explicitly acknowledges these limitations and discusses the relative merits of our approach versus computational fluid dynamics models (p.\,\pageref{chg:R2.27}; see also our response to \ref{r:2.27}).

Regarding the suggestion to model infections under various mixing assumptions: our two model configurations---the proximity-based (non-uniform mixed) model and the well-mixed (uniform) model---effectively represent the two extremes of the mixing spectrum. The proximity-based model assumes radial diffusion without advective mixing, while the well-mixed model assumes complete and instantaneous mixing. The intervention effect is estimated as the ratio between these two models, which inherently captures the range of outcomes under different mixing assumptions. Intermediate mixing scenarios would yield intervention effects that fall between these two bounds. We agree that empirical airflow data would help determine where on this spectrum the true situation lies, and we have noted this as a direction for future research.

We would also like to highlight that our modelling framework is comprehensive in its treatment of uncertainty. We performed 100 Monte Carlo simulations per study day, sampling from prior distributions for quanta generation rates, bacterial inactivation, gravitational settling, and the proportion of presumptive TB patients who are infectious. The resulting credible intervals reflect the cumulative uncertainty across these parameters. We also conducted sensitivity analyses varying assumptions about the relative infectiousness of confirmed vs.\ presumptive TB patients and of individuals with HIV. While there is always scope for further refinement---and empirical validation of modelled results remains an important goal for future studies---we believe our approach provides a useful assessment of the \textit{potential} benefits of proximity-targeting interventions.

We agree with the reviewer that the original claim of ``a significant public health impact'' was too strong given that these are modelled, not empirical, results. We have attenuated the language throughout the manuscript---for instance, replacing ``present a significant public health impact'' with ``may have a meaningful public health impact''---to make clear that these are estimated potential benefits, not direct evidence of prevented infections (see \ref{r:2.1} for a detailed list of changes).}

\item\label{r:2.8} \comment{I am not convinced that the authors' attempt to use clinical data to identify people with TB is a good idea. In outpatient settings, WHO symptom screens will miss most people with TB (see e.g.\ PMID 34864910). This probably explains why the estimate of TB prevalence in clinic attendees (713 / 178510 = 0.4\%) is lower than that estimated in recent community TB prevalence surveys in urban Zambia, e.g.\ TREATS. The WHO symptom screen also has poor specificity. Including people with cough/fever but negative tests for TB as `presumptive cases' sets up secular trends in assumed TB prevalence, likely driven by variation in circulation of respiratory viruses. This is problematic when making a before/after comparison. I think it would be cleaner to fix TB prevalence, using a variety of different assumptions, then to model the impact of their intervention to reduce clustering, fitting to the data on patient numbers, patient locations, and CO2, all of which feel more robust.}

\response{We appreciate the reviewer's thoughtful concerns about using clinical data to identify TB patients. We agree that WHO symptom screens have limited sensitivity and specificity in outpatient settings, and that the resulting estimate of TB prevalence among clinic attendees may underestimate the true burden. However, the modelled intervention effects are not sensitive to variation in the number or identification of TB patients, as we explain below.

Our primary outcome---the intervention effect on Mtb transmission risk---is computed as the risk ratio between a proximity-based model (incorporating spatial positions) and a well-mixed model (ignoring spatial positions). If more or fewer TB patients attend on a given day, this affects both the numerator and the denominator of this ratio in the same way, because both models use the same TB patient data as input. Consequently, secular trends in assumed TB prevalence---whether driven by seasonal variation in respiratory viruses or by misclassification of presumptive cases---do not bias the estimated intervention effects. Using a fixed TB prevalence, as the reviewer suggests, may provide a better approximation of the \textit{absolute} Mtb transmission risk, but it would not change the \textit{relative} effect estimates that are the focus of our analysis.

We also note a practical advantage of using clinical data to identify TB patients rather than assigning TB status at random based on a fixed prevalence. For each clinically identified TB patient, we made considerable effort to reconstruct their complete movement trajectory through manual re-linking of interrupted tracking segments (described in the supplementary material). This ensures that the spatial exposure patterns attributed to TB patients are representative in terms of their locations visited and duration of stay. If TB status were instead assigned randomly to tracked individuals, short or interrupted person movements could be attributed to TB patients, potentially introducing bias into the modelled spatial exposure and transmission risk.

Regarding TB prevalence: the reviewer notes that the estimated TB prevalence among clinic attendees (713/178,510 = 0.4\%) is lower than community-based estimates. We note that the denominator (178,510 entries) overestimates the number of unique clinic attendees, as it includes re-entries and very short entries by non-attending visitors. The TREATS study reported a TB prevalence of 0.51\% for Zambian communities \citep{klinkenberg2023PLoSMed}. When considering both confirmed and presumptive TB patients together, our estimate of 0.4\% is thus reasonably close to this community estimate.

To address the reviewer's concern directly, we have added a sensitivity analysis in which all confirmed and presumptive TB patients are considered infectious---rather than only a random subset of presumptive cases as in the main analysis (p.\,\pageref{chg:R2.8}). This scenario effectively assumes a higher TB prevalence consistent with community-based estimates, while retaining the advantage of using clinically identified patients with well-reconstructed movement trajectories. As expected, this scenario yields a substantially higher absolute modelled transmission risk: the expected number of infections was 172\% (95\% quantile 54--320) higher than in the main analysis. The relative intervention effects, however, remain consistent, as both the proximity-based and well-mixed models are equally affected by the increased number of infectious individuals. The results are reported alongside the other sensitivity analyses in the Results section (p.\,\pageref{chg:R2.8results}) and Supplementary Figure \ref{fig:S7}.}

\item\label{r:2.9} \comment{In presenting the modelled impact on Mtb transmission (abstract and results), it is unclear why the authors sum over both interventions, given they have data that might allow them to look separately at each intervention. Also, surely we want rates, or differences per day or per visit, rather than counts.}

\response{We agree that presenting intervention effects separately is important. To clarify: we do report them individually---the relative reductions in modelled Mtb transmission risk are presented separately for each intervention phase with means and 95\% credible intervals (e.g., 39\% reduction for the first intervention, 21\% for the second). The summed count of prevented infections over the entire intervention period was intended to convey the overall magnitude of the combined effect, given that the two interventions are cumulative (the second retained all measures from the first and added a patient flow system).

We agree with the reviewer that the absolute count of prevented infections requires context and is difficult to interpret without a denominator. We have therefore revised the Abstract and Results to report prevented infections as a daily rate and as a rate per entry ($\approx$\,attendee), in addition to the relative risk reductions (Results: p.\,\pageref{chg:R2.9rates}).}

\item\label{r:2.10} \comment{The authors claim that `individuals with TB-like symptoms can generate aerosolised Mtb even in the absence of a confirmed diagnosis' (methods) is clearly contingent on them having TB, even if this is not detected by standard tests.}

\response{The reviewer is correct---aerosolised Mtb can only be generated by individuals who actually have TB, regardless of whether this is detected by standard diagnostic tests. Our original formulation was imprecise. We have revised the sentence in the Methods to make this contingency explicit: presumptive TB patients are considered potentially infectious because a proportion may have undetected TB, not because symptoms themselves generate Mtb (p.\,\pageref{chg:R2.10}).}

\item\label{r:2.11} \comment{What was assumed about CO2 levels in the replacement air? Is this plausible, i.e.\ was the replacement air coming mostly from outside or was it coming from other occupied clinical spaces? Were there any environmental sources of CO2? In my experience, outdoor CO2 can vary throughout the day -- I suspect mostly related to traffic fumes. Were any measurements of CO2 in the replacement air taken? Given indoor CO2 levels were low, assumptions about CO2 concentration in the replacement air will have a big impact on estimates of ventilation rates.}

\response{Thank you for raising these important points about CO\textsubscript{2} assumptions. We used two complementary approaches to estimate the air exchange rate, both described in the supplementary material: a transient mass balance model and a steady-state model.

In the transient mass balance model, the outdoor CO\textsubscript{2} level is estimated as a model parameter but constrained to the range of 300--500\,ppm. In the steady-state model, the outdoor CO\textsubscript{2} level was set to the minimum indoor level observed during nighttime hours (when the waiting area was usually unoccupied). Reassuringly, the two approaches yielded similar outdoor CO\textsubscript{2} estimates, typically between 300 and 400\,ppm, with the transient model estimate usually being slightly lower than the minimum indoor level. The main analysis reports air exchange rates from the transient mass balance model, and these rates are also used for the Mtb transmission modelling. We have clarified in the manuscript (p.\,\pageref{chg:R2.11envmon}) and Supplementary Text B which outdoor CO\textsubscript{2} level was assumed in each approach and how the two estimates compare.

We did not measure outdoor CO\textsubscript{2} levels directly, which we acknowledge as a limitation. It is possible that traffic-related emissions could elevate outdoor CO\textsubscript{2} during daytime hours, which would mean that using a fixed nighttime-derived baseline could overestimate the indoor--outdoor CO\textsubscript{2} differential and thereby bias the estimated air exchange rates downwards. However, the waiting area was naturally ventilated---primarily through the main entrance door and openable windows, with airflow occasionally supported by ceiling fans---and the hospital is situated in a residential area without nearby combustion sources. Replacement air was therefore predominantly outdoor air, making the outdoor CO\textsubscript{2} concentration the appropriate reference (p.\,\pageref{chg:R2.3vent}).

We have made several clarifications in the revised manuscript in response to this and related comments:

\begin{enumerate}
\item We expanded the Methods---Environmental monitoring subsection to summarise both CO\textsubscript{2}-based estimation approaches (transient mass balance and steady-state models) and their consistency (p.\,\pageref{chg:R2.11envmon}).
\item We updated Supplementary Text B to explicitly compare the air exchange rates estimated by both approaches and their respective assumptions about outdoor (replacement air) CO\textsubscript{2} concentrations, showing that the transient model yields slightly higher air exchange rates and lower outdoor CO\textsubscript{2} estimates than the steady-state model, but that both are consistent.
\item We added a description of the waiting area dimensions (approximately 393\,m\textsuperscript{2}, volume 1,179\,m\textsuperscript{3}) to the Methods---Study setting subsection (p.\,\pageref{chg:R2.11setting}).
\item We added a paragraph to the Environmental monitoring subsection describing the ventilation conditions, airflow sources (natural wind, occasional ceiling fans), negligible air exchange with adjacent clinic areas, the absence of nearby combustion sources, and the predominantly outdoor origin of replacement air (p.\,\pageref{chg:R2.3vent}).
\item We now report absolute ventilation rates (L/s) alongside air exchange rates in the Results section (p.\,\pageref{chg:R2.14}) and Table~\ref{tab:2} (p.\,\pageref{chg:R2.14tab}); see also our response to \ref{r:2.14}.
\item The Methods section has been restructured: ``Study setting and interventions'' has been split into separate ``Study setting'' and ``Study interventions'' subsections (p.\,\pageref{chg:R2.12}), and the former Supplementary Table~S1 describing the interventions has been moved into the main manuscript as Table~\ref{tab:1} (p.\,\pageref{tab:1}).
\item We added a paragraph to the Discussion acknowledging the limitations of not measuring outdoor CO\textsubscript{2} and airflow patterns, and discussing the relative merits of CFD-based and Wells-Riley-based approaches (p.\,\pageref{chg:R2.27}; see also \ref{r:2.27}).
\end{enumerate}}

\item\label{r:2.12} \comment{I'd like to see more detailed description of both the setting and the interventions in the main manuscript. How big is the waiting area? How long do people wait? What is the mix of walk-in patients vs booked appointments? How many people are attending for elective vs emergency care? Cross ventilation or stack ventilation? Why were patients not asked to wait outside?}

\response{We have restructured the Methods section to provide more detailed and clearly separated descriptions of the setting and the interventions. The former ``Study setting and interventions'' subsection has been split into two distinct subsections: ``Study setting'' and ``Study interventions'' (p.\,\pageref{chg:R2.12}). In addition, we moved the former Supplementary Table~S1 into the main manuscript as Table~\ref{tab:1} (p.\,\pageref{tab:1}), so that the full description of each intervention phase is now part of the main text.

Regarding the specific questions raised: the hospital is a busy general facility serving exclusively walk-in patients, with a mix of emergency and elective cases (i.e., there are no booked appointments). We have added this information to the Study setting subsection (p.\,\pageref{chg:R2.12setting}). We did not collect data on individual waiting times, but the distribution of track durations provides a proxy: the median duration per tracked movement was 1.3 minutes (IQR 0.5--4.7), though longer tracks exceeding 30 minutes---likely representing people seated while waiting---accounted for a substantial proportion of total person-time (Supplementary Figure \ref{fig:S3}). Without dedicated crowd management, attendees typically sat as close to the doctors' offices as possible, resulting in concentrated clustering in the front rows. Patients were not routinely asked to wait outside because, absent dedicated staff, attendees preferred to remain near the consultation rooms to avoid missing their turn. During the intervention phases, patients were encouraged to wait outside or to spread more evenly across the waiting area, but sustaining such behaviour required dedicated auxiliary staff---which is part of why the interventions were effective only when actively managed. The waiting area dimensions (approximately 393\,m\textsuperscript{2}), ventilation type (natural cross ventilation through the main entrance door and windows on opposite sides of the building, occasionally supported by ceiling fans), and the absence of mechanical ventilation or air exchange with adjacent clinical spaces are now described in the Study setting (p.\,\pageref{chg:R2.11setting}) and Environmental monitoring (p.\,\pageref{chg:R2.3vent}) subsections.}

\item\label{r:2.13} \comment{I would also like to see more detailed description of the interventions in the main manuscript. All the detail in Table S1 belongs in the main manuscript. In Table S1, the `restriction on the number of people allowed to accompany the patient' is not mentioned in the first intervention but it is implied that this was a component of both interventions. Did many people attend with friends or family? Did the intervention seek to separate relatives? Also, presumably, most time was spent sat still. That being the case, why would we expect a one-way system to work?}

\response{We have addressed this comment in several ways. First, we moved the former Supplementary Table S1 into the main manuscript as Table~\ref{tab:1} (p.\,\pageref{tab:1}), so that the detailed description of all study interventions is now directly available to the reader.

Second, the reviewer correctly noted that the restriction on accompanying persons was only mentioned for the second intervention. In fact, this measure was part of both intervention phases. We have clarified this by adding the restriction to the description of the first intervention in Table~\ref{tab:1} (p.\,\pageref{chg:R2.13a}).

Regarding the accompanying persons: as is common in the African setting, patients were frequently accompanied by friends or relatives, which further increased crowding. We have added this to the Study setting (p.\,\pageref{chg:R2.13relatives}). The hospital management tried to minimise the number of accompanying persons per patient, but this was difficult to enforce and socially not well accepted. The interventions did not specifically seek to separate relatives from patients, but the restriction on accompanying persons (applied during both intervention phases) aimed to reduce the overall number of people in the waiting area.

Regarding the one-way system: while most time in the waiting area was indeed spent seated, the one-way patient flow system was not primarily intended to reduce time spent in proximity while seated, but rather to reduce the number of random close-contact encounters that occur as people move through the waiting area---for example, when arriving, leaving, or moving between service points. We have clarified this aim in the Study interventions subsection (p.\,\pageref{chg:R2.13oneway}).}

\item\label{r:2.14} \comment{I would encourage the authors to report absolute ventilation rates in, e.g., m\textsuperscript{2}/hr rather than ACH. It is the absolute ventilation rate that predicts Mtb transmission risk. In a small room, a much higher number of ACH is required to obtain a safe level of ventilation than in a larger room. This is why target ventilation rates in ACH are unhelpful.}

\response{We agree that absolute ventilation rates are more informative than air changes per hour (ACH) for assessing Mtb transmission risk, as ACH values depend on room volume and are not directly comparable across settings of different sizes. We have added the absolute ventilation rate in L/s alongside the air exchange rate in both the Results text (p.\,\pageref{chg:R2.14}) and Table \ref{tab:2} (p.\,\pageref{chg:R2.14tab}). The median absolute ventilation rate was 2,621\,L/s (IQR 2,044--3,541), reflecting the large and well-ventilated waiting area (approximately 393\,m\textsuperscript{2} floor area, estimated volume 1,179\,m\textsuperscript{3}). We retained the ACH values as well, since they are used as input parameters in our extended Wells-Riley model.}

\item\label{r:2.15} \comment{Throughout, greater clarity about units is needed, e.g.\ are the reductions in spatial clustering and modelled Mtb transmission reported in the abstract absolute or relative differences? What are the units of measurement? E.g.\ is this modelled risk of Mtb infection per person per visit? Elsewhere in the manuscript, rates are mentioned, which may make more sense given healthcare workers presumably have much longer exposure.}

\response{We have clarified the units and the distinction between absolute and relative measures throughout the manuscript. The reductions in spatial clustering (Gini coefficient) and modelled Mtb transmission risk reported in the Abstract and Results are relative reductions, expressed as percentages with 95\% credible intervals. We have made this explicit in the revised text. In addition, we now report TB prevalence among clinic attendees (p.\,\pageref{chg:R2.15prev}) and the modelled number of infections not only as cumulative counts but also as a daily rate and as an incidence per entry (p.\,\pageref{chg:R2.15rate}), providing more interpretable measures that account for differences in exposure duration. We also report the prevented infections as daily rates and per-entry rates by intervention phase (p.\,\pageref{chg:R2.9rates}; see also \ref{r:2.9}).}

\item\label{r:2.16} \comment{It is a little unclear to me what purpose the bioaerosol sampling is serving here. The authors make no comment on sensitivity or specificity. We already know, e.g.\ from infection incidence in healthcare workers, that there is aerosolised Mtb in healthcare facilities in Southern Africa (PMID 25946353).}

\response{The primary purpose of the bioaerosol sampling was to empirically confirm the presence of airborne \emph{Mtb} in the hospital's waiting area, thereby demonstrating that the modelled risk of transmission is grounded in real evidence and not purely hypothetical. While previous studies have documented aerosolised \emph{Mtb} in healthcare facilities in Southern Africa, our study is set in a Zambian hospital with different infrastructure and ventilation characteristics, and direct evidence of airborne \emph{Mtb} in this specific setting had not been established.

Regarding sensitivity and specificity: detection of airborne \emph{Mtb} at the low concentrations typically found in well-ventilated spaces remains technically challenging, and bioaerosol concentrations are known to be highly variable \citep{banholzer2025IntJInfectDis}. False negatives are therefore more likely than false positives---meaning that our rare detections likely underestimate the true frequency of airborne \emph{Mtb} presence. Furthermore, the GeneXpert Ultra assay detects \emph{Mtb} genomic DNA and cannot distinguish between viable and non-viable bacilli. We have added a brief note on these detection limitations to the Methods (p.\,\pageref{chg:R2.16}) and already acknowledge them in the Discussion.}

\item\label{r:2.17} \comment{Table 1 -- it needs to be clearer that the clinical data only relate to people with `presumptive' or `confirmed' TB. Are any such data available for the general clinic population?}

\response{We have revised the caption of Table~\ref{tab:2} (p.\,\pageref{chg:R2.17}) to make explicit that the clinical data (patient characteristics and TB status) relate exclusively to people with presumptive or confirmed TB, not to the general clinic population. No individual-level clinical data were collected for other attendees, as only people flagged through symptom-based screening or with a known TB diagnosis had their demographic and clinical characteristics recorded. The environmental and person-tracking data, by contrast, cover all attendees and the entire waiting area.}

\item\label{r:2.18} \comment{Figure 1 -- it would be good to have a better-quality image. Where are the windows? Where is the replacement air coming from? Could we have a graphical representation of the seating and what the interventions looked like?}

\response{We have updated Figure \ref{fig:1} to provide additional information about the waiting hall, including the positions of benches and windows. Airflow directions were not indicated because airflow patterns were not measured and varied continuously with outdoor wind conditions and door/window openings (see our response to \ref{r:2.3}). We have also added a new Figure \ref{fig:interventions} that illustrates the study setting and seating layout during each of the three study phases (baseline, first intervention, second intervention), providing a graphical representation of what the interventions looked like in practice.}

\item\label{r:2.19} \comment{Figure 3 -- in panel B, the authors could use more distinct colours. In Panel C, why is this data from a single CO2 meter? I remain unclear why the authors only used data from one of the two CO2 meters. Why not pool CO2 measurements to better represent what is happening in the waiting room? This approach has been taken elsewhere (PMID 36962521).}

\response{Regarding panel B (now Figure \ref{fig:3}B due to the addition of the new study interventions figure): the colour scheme was deliberately chosen to encode meaningful groupings. The figure displays nine elements---the overall median (black), the interquartile range (grey), and seven weekday-specific averages---making it inherently difficult to find fully distinct colours. Weekend days share a similar bluish hue to visually distinguish them from weekdays, which use a reddish gradient. We believe this grouping aids interpretation. That said, we have made slight modifications to improve distinguishability: the overall uncertainty band is now lighter, the weekday lines are bolder, and we chose slightly darker colours throughout.

Regarding the CO\textsubscript{2} data: we deliberately chose to use data from the single monitor positioned in the more crowded part of the waiting area (AQM1), rather than pooling or averaging measurements from both monitors. Our rationale was to provide a conservative estimate of ventilation conditions: AQM1 recorded consistently higher CO\textsubscript{2} levels than AQM2 (Supplementary Figure \ref{fig:S2}), and using the higher values yields a lower estimated air exchange rate, which is the more cautious assumption for Mtb transmission modelling.

Averaging the two monitors would give equal weight to both locations, effectively diluting the CO\textsubscript{2} signal from the area where most patients---and thus most potential Mtb exposure---were concentrated. This would produce a higher estimated ventilation rate and thus a lower modelled transmission risk, which we considered less conservative.

The study referenced by the reviewer \citep{beckwith2022PLOSGlobalPublicHealth} pooled CO\textsubscript{2} measurements from multiple monitors in a controlled tracer gas experiment with artificial CO\textsubscript{2} release and decay, where pooling serves to reduce measurement noise under standardised conditions. Our setting differs in that we performed passive, longitudinal monitoring of naturally generated CO\textsubscript{2} over 52 days in a large, heterogeneous space with spatially varying occupancy. In this context, we believe that using the monitor closest to the area of highest patient density is more appropriate than pooling across locations with different exposure characteristics. We have clarified this rationale in the revised manuscript (p.\,\pageref{chg:R2.19}).}

\item\label{r:2.20} \comment{Figure 4 -- given the second intervention included the components of the first intervention, presumably the differences between `First intervention' and `Second intervention' are driven by fewer people with `presumptive TB' (many of whom will have had viral respiratory illnesses)? See also comment \#8.}

\response{We believe there may be a misunderstanding here. The reviewer suggests that the smaller effect during the second intervention may be driven by fewer presumptive TB patients attending during that phase. However, our intervention effects are not subject to this bias, because they are estimated as the ratio between the proximity-based and the well-mixed model (see also our response to \ref{r:2.8}). If fewer TB patients attended during the second intervention, this affects both models equally---fewer quanta are generated in both the proximity-based and well-mixed configurations---and the ratio between them is unchanged. The intervention effects therefore reflect changes in the spatial clustering of people, not in the number of infectious individuals present.

Regarding why the second intervention was less effective than the first: we are also uncertain about the precise reasons and can only speculate, which we do in the Discussion. Our local study team reported difficulties in enforcing the one-way movement pattern, particularly on busy days when the waiting area was crowded. The one-way system may have inadvertently increased local clustering on those days by restricting access to certain areas. In contrast, the first intervention---which redesigned the waiting area layout and introduced physical distancing cues---achieved a more uniform spatial distribution through structural changes that did not depend on active compliance. We discuss these observations in the manuscript (p.\,\pageref{disc:interventions}).}

\item\label{r:2.21} \comment{Text A -- the supplementary material seems to launch into detail very quickly. Some preamble and a general description of the approach -- more detailed than that in Methods -- would be helpful.}

\response{We agree that Text A launched into technical detail too quickly. An introductory paragraph has been added at the beginning of Text A that provides a general description of the approach: it explains that raw tracking data contained interruptions (due to occlusion, sensor handoff, or movement outside the monitored area) and that the tracking system did not identify TB patients, motivating the two post-processing steps---automatic rejoining of interrupted tracks and manual identification of TB patient movements---before describing each in detail. In addition, a general preamble at the beginning of the Supplementary Information now summarises the structure and purpose of each section (Texts A--C and supplementary results).}

\item\label{r:2.22} \comment{Text C -- the authors note that they included a gravitational settling rate. Do the data from Wells, Lurie and others not suggest a) that large droplets make little contribution to Mtb transmission, as they don't reach the lower airway and b) that droplets rapidly attain a mass where they remain suspended on air currents until either inhaled or vented out of the room. If the authors are modelling a contribution to transmission from large droplets, they may be overestimating the likelihood of close contact transmission. See, for example, PMID 18921435.}

\response{Thank you for raising this---there appears to be a misunderstanding that we can clarify. The gravitational settling rate in our model does not represent transmission via large droplets. Rather, it is one of three mechanisms by which infectious quanta are \textit{removed} from the air---alongside ventilation (air exchange) and bacterial inactivation. Gravitational settling applies to the small particle nuclei (1--7\,$\mu$m) that carry \emph{Mtb} and are the relevant size fraction for airborne transmission. Even these small particles settle over time, albeit slowly, and including this removal pathway is standard in aerosol transport models.

We agree with the reviewer that large droplets make little contribution to \emph{Mtb} transmission, as they do not reach the lower airways and settle rapidly. Our model does not simulate large-droplet transmission. The quanta in our model represent infectious doses carried by small aerosol particles that remain suspended and can be inhaled. Gravitational settling simply provides an additional, small removal term for these particles.

The contribution of gravitational settling to overall quanta removal is minor compared to ventilation. In our study, the median air exchange rate was 8\,h$^{-1}$, while the median gravitational settling rate was approximately 1.5\,h$^{-1}$ and the bacterial inactivation rate approximately 1\,h$^{-1}$ \citep{banholzer2025PLOSComputBiol}. Ventilation thus dominated quanta removal by roughly five-fold. The inclusion of gravitational settling follows the approach in our previous model development paper \citep{banholzer2025PLOSComputBiol}, where it was added in response to peer review to provide a more complete representation of quanta removal. Omitting it would marginally increase the modelled quanta concentration but would not materially change the results, particularly the relative intervention effects.}

\item\label{r:2.23} \comment{Text C -- it would be good to see the distribution from which the quanta generation rate was sampled. Did this reflect the marked heterogeneity which is known to be a key feature of TB epidemiology. For example, in Rod Escombe's study there was a patient whose quanta production rate was estimated at 226 q/h (PMID 18798687).}

\response{We have expanded the description of the quanta generation rate distribution in Supplementary Text C. The quanta generation rate is sampled from activity-specific lognormal prior distributions with median rates of 1.08 quanta/h while sitting and 2.81 quanta/h while walking. The lognormal distribution has a long right tail: the 95\%-credible interval extends to 386 quanta/h (sitting) and 987 quanta/h (walking), which readily encompasses the extreme quanta generation rates reported by \citet{escombe2007ClinInfectDis} (226 q/h). This parametrisation was specifically chosen to capture the marked heterogeneity in infectiousness that characterises TB epidemiology, where a small number of ``superspreaders'' may generate the majority of secondary infections. The full distribution is shown in Figure E of the Supplementary Material of our model development paper \citep{banholzer2025PLOSComputBiol}. We have added an explicit reference to this figure and the distributional parameters in the revised Supplementary Text C (p.\,\pageref{chg:R2.23}).}

\item\label{r:2.24} \comment{Table S2 -- this seems to me to constitute an unnecessary risk of inadvertent disclosure of patient data. E.g.\ if someone knew that their friend, aged 36, was diagnosed with TB at this facility on 9 July, they can look up their self reported HIV status. If a table like this is felt to be needed, dates could be converted into e.g.\ day 1, day 7, and ages converted into ranges, e.g.\ 30--39 years.}

\response{We agree with the reviewer that the individual-level linelist data posed an unnecessary risk of inadvertent disclosure of patient information. We have removed Table S2 from the supplementary material. Summary statistics of the clinical data (age, sex, HIV status, TB status, and diagnostic results) are reported in Table~\ref{tab:2} in the main manuscript, which provides sufficient information without identifying individual patients. The underlying anonymised linelist data remain available from the authors upon reasonable request, subject to institutional data governance approval.}

\item\label{r:2.25} \comment{Figure S3 -- I am a bit unclear what this is. Is this how long people spent in a single raster? Are the >30 min `movements' people sat still, or people continuously moving for more than 30 mins? If someone pauses transiently, is that one or two movements?}

\response{Figure \ref{fig:S3} shows a descriptive summary of how long individual person tracks lasted. A ``person movement'' is defined as a single continuous track recorded by the optical sensors, from first to last detection. This is not the time spent in a single raster cell, but the total duration of one track, which may span multiple cells. Shorter tracks (e.g., \textless1 min) typically represent people walking through the waiting area, while longer tracks (\textgreater30 min) represent people sitting in one location for an extended period. If someone pauses transiently, they remain part of the same track as long as the sensor continues to detect them. A track ends when the sensor loses the person---e.g., because they leave the monitored area, enter a care room, or are temporarily occluded---at which point a new track begins if they are re-detected. The automatic rejoining algorithm described in Supplementary Text A was used to re-link such interrupted tracks where possible. We have clarified the caption of Figure \ref{fig:S3} to make the definition of a person movement more explicit.}

\item\label{r:2.26} \comment{Figure 10 -- it seems the use of the ratio hasn't removed the weekday vs weekend differences, presumably because the intervention is more effective at separating people when the waiting room is quieter?}

\response{The reviewer is correct that weekend differences remain visible in Supplementary Figure \ref{fig:S10}, but this is expected and does not indicate a problem with our approach. To explain:

The rate ratio (proximity-based infections divided by well-mixed infections) captures the contribution of close-proximity contact to transmission risk. Dividing by person-time was introduced to remove the gradual temporal trend caused by declining attendance over the study period. This works because attendance and the rate ratio co-move in the same direction over time---both decline gradually as fewer people visit the clinic---so dividing one by the other effectively cancels out this shared trend.

The weekend pattern, however, behaves differently. On weekends, attendance drops sharply, but the rate ratio actually \textit{increases}. This is because the few people who do attend on quieter days still cluster in the same preferred areas near the doctors' offices---there is no pressure to spread out to more distant seats. As a result, the well-mixed model (which assumes uniform dilution across the room) benefits more from the reduced attendance than the proximity-based model (which captures that the people present are still clustered near the source). The ratio therefore rises on weekends. Since person-time and the rate ratio move in opposite directions on weekends, dividing by person-time amplifies rather than removes this pattern.

This weekend effect is unrelated to the intervention effects, which are our primary interest. The weekday fixed effects in our regression model absorb these day-of-week differences, ensuring that the estimated intervention effects are not confounded by the recurring weekend pattern. The weekday coefficients themselves are nuisance parameters of no further substantive interest. We have added a remark in the Results section explaining this pattern and the rationale for the weekday adjustment (p.\,\pageref{chg:R2.26}).}

\item\label{r:2.27} \comment{Would the questions the authors propose not be better answered using e.g.\ computational fluid dynamics modelling to simulate the mixing of air, which seems central to understanding whether improved separation of individuals within a single space is likely to have an impact on Mtb transmission? It would be good to see greater discussion about the pros and cons of various modelling approaches.}

\response{We agree that a more detailed discussion of the modelling approach and its alternatives is warranted, and we have added a dedicated paragraph to the Discussion (p.\,\pageref{chg:R2.27}).

Computational fluid dynamics (CFD) models can simulate detailed airflow patterns and their influence on aerosol transport, and would in principle provide a more accurate representation of how infectious quanta are redistributed within the waiting area. However, CFD models have important practical limitations that made them unsuitable for our study. First, they are computationally expensive: simulating airflow in a space of approximately 393\,m\textsuperscript{2} at the temporal resolution required to track dynamic occupancy patterns over 52 study days would be prohibitive. Second, CFD models are primarily designed for experimental or static environments with well-defined boundary conditions (fixed geometry, steady-state airflow). In our setting, airflow conditions changed continuously due to varying wind, door and window openings, occupancy levels, and the occasional use of ceiling fans---making it difficult to specify realistic boundary conditions for a longitudinal study. Third, our modelling approach required many Monte Carlo simulations (100 per study day) to propagate uncertainty in quanta generation, diffusion, and removal parameters, which demands computational efficiency that CFD cannot provide.

Our spatiotemporal extension of the Wells-Riley model was chosen precisely because it balances biological realism with computational feasibility. It incorporates person movements and proximity-dependent exposure---going beyond the classical well-mixed assumption---while remaining tractable for the volume of longitudinal data in our study. We acknowledge that the model does not capture directional airflow or localised air currents, which may affect the spatial quanta concentration. However, since intervention effects are estimated in relative terms (proximity-based vs.\ well-mixed model), systematic biases from the well-mixed assumption would affect both models similarly and are unlikely to qualitatively change the estimated intervention effects.

We note that emerging rapid assessment methods that combine behavioural data with simplified airflow assumptions may help bridge the gap between computationally intensive CFD models and the Wells-Riley framework in future studies. We have cited recent work in this direction and highlighted the incorporation of airflow data as an important direction for future research (p.\,\pageref{chg:R2.27}).}

\end{enumerate}

\subsection*{Minor Comments}

\begin{enumerate}[label=\textbf{R2.\arabic*.}, leftmargin=*, start=28]

\item\label{r:2.28} \comment{Please include page and line numbers.}

\response{We apologise for not including these in the original submission. Page and line numbers have been added to the revised manuscript and supplementary material.}

\item\label{r:2.29} \comment{Please ensure that the distinction between TB (the disease) and Mtb (the pathogen) is maintained throughout. There is no such thing as `TB transmission'.}

\response{We agree---strictly speaking, it is \emph{Mtb} (the pathogen) that is transmitted, not TB (the disease), and we thank the reviewer for this correction. We have reviewed the entire manuscript and supplementary material to ensure that TB refers to the disease and \emph{Mtb} to the pathogen, and that transmission is consistently attributed to \emph{Mtb}. One instance of ``TB transmission'' was corrected to ``\emph{Mtb} transmission'' (p.\,\pageref{chg:R2.29}). The remaining uses of ``TB'' in the manuscript refer appropriately to the disease (e.g., TB case definitions, TB patients, TB burden, TB infection control) rather than to pathogen transmission. We have also updated the title from ``tuberculosis transmission'' to ``\emph{Mycobacterium tuberculosis} transmission''.}

\item\label{r:2.30} \comment{The authors claim that their intervention is low-cost (abstract). In reality, sustaining changes in patient flow would likely require dedicated staff time (PMID 36792227).}

\response{We agree that the term ``low-cost'' was misleading, as sustaining changes in patient flow and spatial management requires dedicated staff time and thus involves continuous costs. We have removed ``low-cost'' from the Abstract (p.\,\pageref{chg:R2.30abs}) and Conclusions (p.\,\pageref{chg:R2.30conc}), replacing it with ``simple'' to more accurately characterise the nature of the interventions. In the Discussion, we have added a caveat noting that sustaining such changes in practice may require dedicated staff time and continuous costs, citing the study referenced by the reviewer (p.\,\pageref{chg:R2.30disc}).}

\item\label{r:2.31} \comment{The scenario where HIV positive people were six times more infectious is not plausible. Yes, in Rod Escombe's animal chamber work, they estimated a quanta generation rate of $\sim$8q/h. However, these were also people with advanced DRTB, which may explain the discrepancy between these findings and those from e.g.\ the Johns Hopkins Pilot Ward. There is a lot of other evidence to suggest that people with HIV associated immunocompromise are less infectious than people not living with HIV (PMID 24368620).}

\response{We thank the reviewer for this correction. We have reviewed the literature and agree that the original assumption of HIV-positive patients being six times more infectious was not supported by the broader evidence. As the reviewer notes, the high quanta generation rates estimated in \citeauthor{escombe2007ClinInfectDis}'s animal chamber work likely reflected advanced drug-resistant TB rather than HIV status per se. We have revised this sensitivity analysis: based on \citet{martinez2021ClinInfectDis}, we now assume that HIV-positive TB patients generate quanta at a 33\% lower rate than HIV-negative patients, reflecting evidence that HIV-associated immunocompromise reduces infectiousness. We have updated the Methods (p.\,\pageref{chg:R2.31}) and re-run the sensitivity analysis accordingly. Under this revised assumption, the expected number of \emph{Mtb} infections was 3\% (95\% quantile --20--0) lower than in the main analysis, consistent with the reviewer's expectation that HIV-associated immunocompromise reduces infectiousness. The intervention effects remain unchanged, as the reduced quanta generation affects both the proximity-based and well-mixed models equally.}

\item\label{r:2.32} \comment{Do the authors expect the intervention to impact waiting time? If so, normalising by person time could remove some of the intervention effect. I expect not, but would be good to include data or a comment.}

\response{We agree with the reviewer's expectation: the interventions were designed to improve the spatial management of attendees already present in the waiting area---reducing crowding and close-proximity contact---not to shorten waiting times. Any effect on waiting time would be incidental. Normalising by person-time therefore does not remove part of the intervention effect. We have added a clarifying statement to the Study interventions subsection (p.\,\pageref{chg:R2.32}).}

\item\label{r:2.33} \comment{There is mention of `masks' being worn by staff members, although this is not modelled. It would be good to be clear whether these are surgical masks, which reduce infectiousness by $\sim$50\% but offer little protection, or N95/FFP3 respirators, which will protect the wearer.}

\response{We have clarified that hospital staff wore surgical masks (p.\,\pageref{chg:R2.33}). Mask use was deliberately not modelled, in order to isolate the impact of the proximity-targeting interventions from other IPC measures already in place. This is particularly relevant when comparing the modelled infection rate between healthcare workers (who consistently wore masks) and other attendees (who rarely did), as reported in the supplementary results---the similar rates between the two groups likely reflect the unmodelled protective effect of surgical masks among staff.}

\item\label{r:2.34} \comment{The authors state that `The number of people in the waiting area peaked at 188\dots' (results) Is this median peak occupancy or the maximum number of people observed in the waiting room on the busiest day?}

\response{This is the median daily peak occupancy (i.e., the median across study days of the maximum number of people present at any point during each day), not the maximum on the single busiest day. We have revised the wording to make this explicit (p.\,\pageref{chg:R2.34}).}

\item\label{r:2.35} \comment{Same paragraph -- not sure we need to be given examples of weekdays!}

\response{We have replaced the examples of weekday names with numerical data: the mean daily peak occupancy per weekday, contrasting Mondays (253) and Tuesdays (232) with Saturdays (144) and Sundays (102), which provides more informative context for the weekday variation (p.\,\pageref{chg:R2.35}).}

\item\label{r:2.36} \comment{Figures S9 and S10 -- again, the authors could chose more distinct colours.}

\response{We have updated the colour scheme in Supplementary Figures \ref{fig:S9} and \ref{fig:S10}, as well as the figure showing the intervention effects (Figure \ref{fig:4}), to use more distinct colours.}

\end{enumerate}

\clearpage

% ============================================================
\section*{Reviewer \#3}
% ============================================================

\noindent\textit{``This is an interesting paper on an important area. The methodological part is generally sound and the statistical methods seem to be applied reasonably well. This is apparent from the supplement which is well written and describes the methods in sufficient detail.''}

\medskip
\noindent We thank Reviewer \#3 for the positive assessment of the methodology and statistical approach. In response to the specific comments, we have clarified the cumulative design of the interventions, the rationale for not including a wash-out period, and the role of the Gini coefficient. All points are addressed below.

\begin{enumerate}[label=\textbf{R3.\arabic*.}, leftmargin=*]

\item\label{r:3.1} \comment{Can the two interventions be assumed to act independently? Are the authors assuming their effect to be additive? Or they are likely to have antagonistic effect if applied simultaneously?}

\response{The two interventions were designed to be cumulative: the second intervention retained all measures from the first (redesigned layout, signage, distancing reminders) and added a one-way patient flow system. They were not assumed to act independently or additively in the statistical model---rather, we estimated the effect of each study phase separately relative to baseline, without decomposing the second intervention into a ``first intervention effect'' plus an ``additional flow system effect.''

In principle, the patient flow system should complement the distancing measures, as directing movement aims to reduce random close-contact encounters independently of seating arrangements. We see no mechanism by which the flow system would be antagonistic to physical distancing. However, in practice, we observed that the second intervention did not produce larger reductions than the first, possibly because the one-way movement pattern was difficult to enforce on busy days and may have inadvertently increased local clustering by restricting access to certain areas. We discuss this in the manuscript (p.\,\pageref{disc:interventions}).}

\item\label{r:3.2} \comment{The question on disentangling the two effects links to the timelines of the two interventions. It seems like the second intervention was essentially applied immediately after the first. Should a ``wash out'' period be considered? Could such a period have disentangled the effect of the two interventions clearly?}

\response{We agree that a wash-out period between the two intervention phases could have provided additional assurance that there was no carryover from the first to the second intervention---for example, if clinic staff remembered procedures from the first phase or if patients retained behavioural changes. However, as discussed in our response to \ref{r:1.2}, we believe such carryover effects were limited: all infrastructure changes were temporary and removed between phases, study staff were explicitly instructed not to continue intervention-related activities outside the designated phases, and re-occurring patients within the short study period were rare (see \ref{r:1.2}, p.\,\pageref{chg:R1.2}). Given these considerations, we believe the absence of a wash-out period is unlikely to have materially affected the results, although incorporating one at the study design stage would be advisable for future studies of this kind.}

\item\label{r:3.3} \comment{The basic modelling assumptions appear reasonable. I wonder if an alternative to the Gini coefficient is worth trying. Kernel smoothing uses many different models, with power or exponential decay over distance being more common. Did the authors investigate the sensitivity of their findings to this choice?}

\response{Thank you for this suggestion. We have added a sensitivity analysis examining both an alternative kernel method and an alternative measure of spatial evenness (p.\,\pageref{chg:R3.3}).

For the kernel, we replaced the Gaussian kernel with an exponential kernel ($\exp(-r/\sigma)$), which has a sharper drop-off and heavier tails than the Gaussian, giving more weight to immediate neighbours while still capturing more distant contributions. The exponential kernel is a natural alternative because it is widely used in spatial epidemiology to model distance-dependent interactions and because its decay profile may better reflect the rapid dilution of infectious aerosols with distance from the source.

For the measure, we computed Shannon entropy as an alternative to the Gini coefficient. While the Gini quantifies inequality in the spatial distribution (higher values = more clustering), entropy quantifies evenness (higher values = more uniform distribution). Both capture spatial clustering from complementary perspectives.

Results were qualitatively similar across all combinations of kernel and measure: the interventions reduced spatial clustering during both intervention phases, with larger reductions during the first intervention (Supplementary Table \ref{tab:S3}, p.\,\pageref{chg:R3.3results}).}

\item\label{r:3.4} \comment{How does the potential for HIV positivity affect the results?}

\response{HIV positivity is considered in our model only through a sensitivity analysis examining its impact on the quanta generation rate. In the original submission, we assumed that HIV-positive TB patients generated quanta at six times the rate of HIV-negative patients, following \citet{escombe2007ClinInfectDis}. In response to Reviewer \#2 (see \ref{r:2.31}), we have revised this scenario: the evidence now suggests that people with HIV-associated immunocompromise are likely \textit{less} infectious, so we have reversed the direction of this sensitivity analysis accordingly. Under the revised assumption, the expected number of \emph{Mtb} infections was 3\% (95\% quantile --20--0) lower than in the main analysis, consistent with reduced infectiousness among HIV-positive TB patients.

We did not model the effect of HIV on susceptibility. Attendees living with HIV but without TB co-infection may have increased susceptibility to \emph{Mtb} infection given the same dose, which could affect the absolute number of modelled infections. We have added a brief note on this in the limitations (p.\,\pageref{chg:R3.4}). However, since the intervention effects are estimated in relative terms (as the ratio between the proximity-based and well-mixed models), differential susceptibility among attendees does not affect the estimated intervention effects.}

\item\label{r:3.5} \comment{As the authors mention, TB transmission is incompletely understood and therefore this is very much an area where ``all models are wrong'' -- a fact which ought to be taken into account when interpreting the paper's findings. The authors added sufficient nuance and interval estimates throughout so that helps with not over-interpreting the paper's findings.}

\response{We thank the reviewer for this balanced perspective and agree that all models are simplifications of reality. We have taken care throughout the manuscript to present credible intervals, sensitivity analyses, and caveats about the modelling assumptions. In response to comments from all reviewers, we have further attenuated the language to make clear that our findings represent modelled estimates of potential benefits rather than empirical evidence of prevented infections (see \ref{r:2.1} and \ref{r:2.7} for details of these changes).}

\end{enumerate}

\clearpage

\section*{Revised and New Figures}

For the convenience of the reviewers, we include below copies of all figures that were revised or newly added.

\begin{figure}[H]
\centering
\includegraphics[width=0.9\textwidth]{../../../revised_study-setting.png}
\caption*{\textbf{Figure \ref{fig:1} (revised).} Study setting, now with positions of benches and windows indicated.}
\end{figure}

\begin{figure}[H]
\centering
\includegraphics[width=0.9\textwidth]{../../../revised_study_phase_baseline.png}\\[1em]
\includegraphics[width=0.9\textwidth]{../../../revised_study_phase_first-intervention.png}\\[1em]
\includegraphics[width=0.9\textwidth]{../../../revised_study_phase_second-intervention.png}
\caption*{\textbf{Figure \ref{fig:interventions} (new).} Illustration of the intended impact of study interventions. (A) Baseline: routine clinic conditions with patients clustering near registration and consultation rooms. (B) First intervention: redesigned seating layout for more uniform spread. (C) Second intervention: added one-way patient flow system.}
\end{figure}

\begin{figure}[H]
\centering
\includegraphics[width=0.9\textwidth]{../../../../results/modelled_data.png}
\caption*{\textbf{Figure \ref{fig:3} (revised).} Clinical, environmental and person-tracking data, with updated colour scheme in panel B.}
\end{figure}

\begin{figure}[H]
\centering
\includegraphics[width=0.9\textwidth]{../../../../results/modelled_effects.png}
\caption*{\textbf{Figure \ref{fig:4} (revised).} Effects of proximity-targeting interventions on spatial clustering and \emph{Mtb} transmission, with updated colour scheme.}
\end{figure}

\bibliographystyle{apalike}
\bibliography{../references}

\end{document}
